\section{A Uniqueness Theorem}

In this subsection, for any strictly positive integer m, any integer i in the interval \([1,m]\) , and any sequence \((v_{1},\ldots,v_{m})\) in \(V^{m}\) , we write \((v_{1},\ldots,\widehat{v_{i}},\ldots,v_{m})\) for the sequence \((v_{1},\ldots,v_{i-1},v_{i+1},\ldots,v_{m})\) in \(V^{m-1}\) .

PROPOSITION 1. --- Let W be a hyperplane in V and e be a vector in V---W. Let \(\varphi\) be an element of \(\Omega(\mathrm{W})\) . There exists a unique element \(\omega\) of \(\Omega(\mathrm{V})\) such that we have

\[
\omega (w _ {1}, \ldots , w _ {n - 1}, e) = \varphi (w _ {1}, \ldots , w _ {n - 1})\tag{7}
\]

for every basis \((w_{1},\ldots ,w_{n - 1})\) of W.

\begin{enumerate}
\def\labelenumi{\Alph{enumi})}
\tightlist
\item
  Uniqueness of \(\omega\) : Let \(\omega\) be an element of \(\Omega(\mathrm{V})\) . Let \((v_{1},\ldots,v_{n})\) be a basis of V, and let \(\mu_{1},\ldots,\mu_{n}\) be the coordinates of e in this basis. Denote by I the set of integers from 1 to n (inclusive) such that \(\mu_{i}\neq0\) ; it is not empty. Let i be an element of I. The sequence \((v_{1},\ldots,\widehat{v_{i}},\ldots,v_{n},e)\) is a basis of V, and since \(e-v_{i}\mu_{i}\) is a linear combination of the sequence \((v_{1},\ldots,\widehat{v_{i}},\ldots,v_{n})\) , formula (6) implies
\end{enumerate}

\[
\omega (v _ {1}, \ldots , \widehat {v _ {i}}, \ldots , v _ {n}, v _ {i} \mu_ {i}) = \omega (v _ {1}, \ldots , \widehat {v _ {i}}, \ldots , v _ {n}, e).\tag{8}
\]

It follows from formulas (1) and (3) that the left-hand side of this equality is equal to \(\pi(-1)^{n-i}\pi(\mu_i)\omega(v_1,\ldots,v_n)\) . Denote by p the projector of V with image W and kernel eD. The vectors \(v_j - p(v_j)\) are proportional to e, and repeatedly applying formula (6) shows that the right-hand side of formula (8) is equal to \(\omega(p(v_1),\ldots,\widehat{p(v_i)},\ldots,p(v_n),e)\) .

It follows from this that if \(\omega\) satisfies relation (7), then we have

\[
\omega (v _ {1}, \dots , v _ {n}) = \pi (- 1) ^ {n - i} \pi (\mu_ {i}) ^ {- 1} \varphi (p (v _ {1}), \dots , \widehat {p (v _ {i})}, \dots , p (v _ {n}))\tag{9}
\]

for every i in I, which proves the uniqueness of \(\omega\) .

\begin{enumerate}
\def\labelenumi{\Alph{enumi})}
\setcounter{enumi}{1}
\tightlist
\item
  Construction of \(\omega\) : Let \((v_{1},\ldots,v_{n})\) be a basis of V; we define \(\mu_{1},\ldots,\mu_{n}\) and I as above. For every \(i\) in I, the assumption \(\mu_{i}\neq0\) implies that the sequence \((v_{1},\ldots,\widehat{v_{i}},\ldots,v_{n},e)\) generates the space V, so the sequence \((p(v_{1}),\ldots,\widehat{p(v_{i})},\ldots,p(v_{n}))\) generates the space W = p(V). In other words, this last sequence is a basis of W, and
\end{enumerate}

\[
t _ {i} = \pi (- 1) ^ {n - i} \pi (\mu_ {i}) ^ {- 1} \varphi (p (v _ {1}), \ldots , \widehat {p (v _ {i})}, \ldots , p (v _ {n}))\tag{10}
\]

defines an element of \(\mathrm{D}_{\mathrm{ab}}^{*}\) .

Let i and j be two elements of I such that i \textless{} j; let us prove the equality \(t_{i} = t_{j}\) . By definition, the vector \(v_{i}\mu_{i} + v_{j}\mu_{j} - e\) is a linear combination of the vectors \(p(v_{k})\) for k different from i and j. Consequently, \(p(v_{i})\mu_{i} + p(v_{j})\mu_{j}\) is a linear combination of the vectors \(p(v_{k})\) for k different from i and j. By formula (6), we therefore have

\[
\begin{array}{c} \varphi (p (v _ {1}), \ldots , \widehat {p (v _ {i})}, \ldots , \widehat {p (v _ {j})}, \ldots , p (v _ {n}), p (v _ {i}) \mu_ {i}) \\ = \varphi (p (v _ {1}), \ldots , \widehat {p (v _ {i})}, \ldots , \widehat {p (v _ {j})}, \ldots , p (v _ {n}), - p (v _ {j}) \mu_ {j}). \end{array}
\]

By applying formulas (1) and (3), we deduce

\[
\begin{array}{c} \varphi (p (v _ {1}), \ldots , \widehat {p (v _ {j})}, \ldots , p (v _ {n})) \pi (\mu_ {i}) \pi (- 1) ^ {n - i - 1} \\ = \varphi (p (v _ {1}), \ldots , \widehat {p (v _ {i})}, \ldots , p (v _ {n})) \pi (- \mu_ {j}) \pi (- 1) ^ {n - j}, \end{array}
\]

and therefore \(t_i = t_j\) .

Having proved this, we define \(\omega(v_1, \ldots, v_n)\) as the common value of the \(t_i\) for \(i\) in I. We have thus constructed a mapping \(\omega\) from B(V) to D \(_{\text{ab}}^*\) that satisfies relation (9). Let \((w_1, \ldots, w_{n-1})\) be a basis of W; if we set \(v_i = w_i\) for \(1 \leqslant i \leqslant n-1\) and \(v_n = e\) , then we have I = \{n\} and \(\mu_n = 1\) . We moreover have \(p(v_i) = w_i\) for \(1 \leqslant i \leqslant n-1\) , and the formula \(\omega(w_1, \ldots, w_{n-1}, e) = \varphi(w_1, \ldots, w_{n-1})\) is a specific case of (9).

We must now prove that \(\omega\) belongs to \(\Omega (\mathrm{V})\) .

\begin{enumerate}
\def\labelenumi{\Alph{enumi})}
\setcounter{enumi}{2}
\tightlist
\item
  Proof of formula (1): Let \(\lambda_1, \ldots, \lambda_n\) be elements of \(D^*\) and \((v_1, \ldots, v_n)\) be a basis of \(V\) . We define \(\mu_1, \ldots, \mu_n\) and \(I\) as above. Choose an \(i\) in \(I\) . Since we have \(e = \sum_{j=1}^{n} (v_j \lambda_j)(\lambda_j^{-1} \mu_j)\) , formula (9) implies
\end{enumerate}

\[
\begin{array}{l l} (1 1) & \omega (v _ {1} \lambda_ {1}, \ldots , v _ {n} \lambda_ {n}) \\ & \qquad = \pi (- 1) ^ {n - i} \pi (\lambda_ {i} ^ {- 1} \mu_ {i}) ^ {- 1} \varphi (p (v _ {1}) \lambda_ {1}, \ldots , \widehat {p (v _ {i}) \lambda_ {i}}, \ldots , p (v _ {n}) \lambda_ {n}). \end{array}
\]

But we have \(\pi (\lambda_i^{-1}\mu_i)^{-1} = \pi (\mu_i)^{-1}\pi (\lambda_i)\) and

\[
\begin{array}{c} \varphi (p (v _ {1}) \lambda_ {1}, \ldots , \widehat {p (v _ {i}) \lambda_ {i}}, \ldots , p (v _ {n}) \lambda_ {n}) \\ = \varphi (p (v _ {1}), \ldots , \widehat {p (v _ {i})}, \ldots , p (v _ {n})) \pi (\lambda_ {1} \ldots \widehat {\lambda_ {i}} \ldots \lambda_ {n}). \end{array}
\]

Comparing formulas (9) and (11) leads to the desired relation

\[
\omega (v _ {1} \lambda_ {1}, \ldots , v _ {n} \lambda_ {n}) = \omega (v _ {1}, \ldots , v _ {n}) \pi (\lambda_ {1} \ldots \lambda_ {n}).
\]

\begin{enumerate}
\def\labelenumi{\Alph{enumi})}
\setcounter{enumi}{3}
\tightlist
\item
  Proof of the formula (2): Let \((v_{1},\ldots,v_{n})\) be a basis of V and i,j be two distinct integer in the interval {[}1,n{]}; we define \(\mu_{1},\ldots,\mu_{n}\) as before. Consider the basis \((v_{1}^{\prime},\ldots,v_{n}^{\prime})\) of V defined by \(v_{i}^{\prime}=v_{i}+v_{j}\) and \(v_{k}^{\prime}=v_{k}\) for \(k\neq i\) , and let us introduce the coordinates \(\mu_{1}^{\prime},\ldots,\mu_{n}^{\prime}\) of e with respect to this basis; they satisfy \(\mu_{j}^{\prime}=\mu_{j}-\mu_{i}\) and \(\mu_{k}^{\prime}=\mu_{k}\) for \(k\neq j\) . We set \(t=\omega(v_{1},\ldots,v_{n})\) and \(t^{\prime}=\omega(v_{1}^{\prime},\ldots,v_{n}^{\prime})\) . We must prove that t and \(t^{\prime}\) are equal.
\end{enumerate}

Let us first note that, by the definition of \(\omega\) , we have

\begin{enumerate}
\def\labelenumi{(\arabic{enumi})}
\setcounter{enumi}{11}
\tightlist
\item
\end{enumerate}

\[
t = \pi (- 1) ^ {n - k} \pi (\mu_ {k}) ^ {- 1} \varphi (p (v _ {1}), \ldots , \widehat {p (v _ {k})}, \ldots , p (v _ {n})),\tag{13}
\]

\[
t ^ {\prime} = \pi (- 1) ^ {n - k} \pi (\mu_ {k} ^ {\prime}) ^ {- 1} \varphi (p (v _ {1} ^ {\prime}), \ldots , \widehat {p (v _ {k} ^ {\prime})}, \ldots , p (v _ {n} ^ {\prime}))
\]

for every \(k\) such that \(\mu_{k}\) and \(\mu_k^\prime\) are nonzero.

\begin{enumerate}
\def\labelenumi{\alph{enumi})}
\item
  If \(\mu_i \neq 0\) , then the sequences \((p(v_1), \ldots, \widehat{p(v_k)}, \ldots, p(v_n))\) and \((p(v_1'), \ldots, \widehat{p(v_k'}, \ldots, p(v_n'))\) are equal, and we have \(\mu_i = \mu_i'\) , and therefore \(t = t'\) .
\item
  If there exists an index \(k\) different from \(i\) and \(j\) such that \(\mu_k \neq 0\) , then we have \(p(v_l') = p(v_l)\) for \(l \neq i\) and \(p(v_i') = p(v_i) + p(v_j)\) . The elements \(p(v_i)\)
\end{enumerate}

and \(p(v_{j})\) both belong to the sequence \((p(v_{1}),\ldots ,\widehat{p(v_{k})},\ldots ,p(v_{n}))\) . Since \(\varphi\) belongs to \(\Omega (\mathrm{W})\) , formula (2) of VIII, p.~447 applies; it gives

\[
\varphi (p (v _ {1}), \dots , \widehat {p (v _ {k})}, \dots , p (v _ {n})) = \varphi (p (v _ {1} ^ {\prime}), \dots , \widehat {p (v _ {k} ^ {\prime})}, \dots , p (v _ {n} ^ {\prime})).
\]

But we also have \(\mu_{k} = \mu_{k}^{\prime}\) , and therefore \(t = t'\) .

\begin{enumerate}
\def\labelenumi{\alph{enumi})}
\setcounter{enumi}{2}
\tightlist
\item
  It remains to examine the case when the only index \(k\) such that \(\mu_k \neq 0\) is \(j\) . We then have \(\mu_j' = \mu_j\) , \(e = v_j\mu_j\) , and \(p(v_j) = 0\) . Set \(k = j\) in formulas (12) and (13). Since the sequences \((p(v_1), \ldots, \widehat{p(v_j)}, \ldots, p(v_n))\) and \((p(v_1'), \ldots, \widehat{p(v_j'), \ldots, p(v_n'))}\) are equal, we have \(t = t'\) .
\end{enumerate}

COROLLARY 1. --- Let \((e_1, \ldots, e_n)\) be a basis of V and t an element of \(\mathrm{D}_{\mathrm{ab}}^*\) . There exists a unique element \(\omega\) of \(\Omega(\mathrm{V})\) such that \(\omega(e_1, \ldots, e_n) = t\) .

Let us prove the corollary by induction on n.~If n = 0, then the only basis of V is the empty one, so the assertion is true. Now suppose \(n \geqslant 1\) ; denote by W the subspace of V generated by \(e_{1}, \ldots, e_{n-1}\) . By Proposition 1, there exists a bijection \(\Lambda\) from \(\Omega(\mathrm{V})\) to \(\Omega(\mathrm{W})\) satisfying

\[
(\Lambda (\omega)) (w _ {1}, \ldots , w _ {n - 1}) = \omega (w _ {1}, \ldots , w _ {n - 1}, e _ {n})
\]

for every basis \((w_{1},\ldots,w_{n-1})\) of W and every \(\omega\) in \(\Omega(\mathrm{V})\) . By the induction hypothesis, there exists a unique element \(\varphi\) of \(\Omega(\mathrm{W})\) such that \(\varphi(e_{1},\ldots,e_{n-1})=t\) . The relation \(\omega(e_{1},\ldots,e_{n})=t\) , for \(\omega\) in \(\Omega(\mathrm{V})\) , is equivalent to \(\Lambda(\omega)=\varphi\) . Corollary 1 follows.

COROLLARY 2. --- Let \(\omega\) and \(\omega'\) be two elements of \(\Omega(V)\) . There exists a unique element \(t\) of \(\mathrm{D}_{\mathrm{ab}}^{*}\) such that \(\omega' = t\omega\) .

Remark. --- Suppose that the field D is commutative. By definition, B(V) is a subset of \(V^{n}\) . It is clear that the restriction to B(V) of a nonzero alternating n-linear form \(f: V^{n} \to D\) belongs to \(\Omega(V)\) . Moreover, if \((e_{1}, \ldots, e_{n})\) is a basis of V and t is a nonzero element of D, then there exists a unique alternating n-linear form f such that \(f(e_{1}, \ldots, e_{n}) = t\) (III, §7, No.~4, p.~511 and III, §7, No.~8, p.~518). By Corollary 1, the set \(\Omega(V)\) consists of the restrictions to B(V) of the nonzero alternating n-linear forms.

